\subsection{One-Shot Deadline Timer}\label{section:deadline-timer}
Each (:term:base.terms.logical_processor:) owns one one-shot deadline timer consisting of armed state, an absolute
timebase deadline, and a 24-bit interrupt identity. TARM and TCANCEL are the only architectural programming
operations. The timer state has no control-register or SAVE-area projection.

TARM observes the current time and classifies its absolute deadline by signed modular distance:

\[
  \mathit{distance}=\operatorname{signed64}(\mathit{deadline}-\mathit{time}).
\]

A positive distance arms the timer and atomically replaces any previous arm. Software can therefore schedule at
most $2^{63}-1$ ticks into the future with one operation. A zero or negative distance is already reached: TARM posts
the supplied identity to the local interrupt file and leaves the timer disarmed. The supplied 64-bit identity value
must be from 1 through ICAP.MAX\_ID. An invalid identity raises
(:ref:base.events.EXCEPTION.INVALID_CONTROL_SELECTOR:) without changing timer state.

When time reaches an armed deadline, posting the configured identity and clearing armed state are one atomic
transition. Expiry is one-shot; periodic behavior requires software to install another absolute deadline. TCANCEL
atomically clears armed state and is a no-op when already disarmed. Neither TARM nor TCANCEL clears a pending bit
created by an earlier expiry.

Expiry posts pending state regardless of the identity\textquotesingle{}s enable bit, ITHRESH, STATUS.IE, ECR.V, or
event depth. Those gates remain owned by the interrupt file and event-admission model. A halted logical processor
leaves HALT only if the resulting interrupt becomes admissible; otherwise the identity remains pending. A WAIT may
return spuriously so that an admissible interrupt can be processed, while time advancement and posting remain
independent of that return.

Time advancement, expiry, TARM, TCANCEL, interrupt-file pending operations, and event-entry commit are observed in
one of their legal serial orders. A cancellation ordered before expiry prevents the post; a cancellation ordered
after expiry cannot remove it. An expiry ordered before replacement may leave the old identity pending, whereas a
replacement ordered first prevents the superseded deadline from posting. No operation loses a post that has already
committed. These operations do not by themselves order memory operations.

Warm and cold reset disarm the timer. Reset does not clear interrupt-file pending state except through the interrupt
file\textquotesingle{}s own reset transition.
